\providecommand{\zkTBC}[1]{\textbf{[TO COMPLETE: #1]}}

\begin{abstract}
Zero-knowledge machine learning (zkML) lets clients verify outsourced inference, but generating proofs is slow and memory intensive. Splitting a model into smaller proof jobs can reduce proving time and peak memory and allow jobs to run in parallel. However, a runtime must still decide how jobs from concurrent requests share the available compute and memory.

We characterize partitioned zkML execution and identify three properties relevant to scheduling: jobs vary widely in runtime and peak memory, a job's memory demand changes across proving stages, and repeated requests can reuse compiled circuits and proving keys. We present \emph{zkInfer}, a distributed runtime that schedules these jobs across prover nodes. A coordinator assigns jobs from active requests to nodes. Each node uses recent memory measurements to decide when a job can begin its next stage, rather than reserving its peak memory for its entire execution. zkInfer also reuses compiled circuits and proving keys without changing the underlying proving protocol.

Stage-level memory admission is most useful when available memory limits concurrency. With MobileNetV2 restricted to a 48~GiB proving budget, zkInfer reduces median request latency by 4.2--8.8\% across bursts of 2--16 requests compared with whole-job peak reservation. Separately, increasing the prover pool from one to ten nodes speeds up a single partitioned request by \(10.5\)--\(12.7\times\) across TinyBERT, MobileNetV2, and NanoGPT. Reusing setup artifacts removes 96.3--98.6\% of repeated setup time. These results show how distributed execution, stage-level memory admission, and setup reuse can improve zkML proof serving.
\end{abstract}
